\documentclass[10pt,a4paper]{amsart}
\usepackage[margin=19mm]{geometry}
\usepackage{amsmath,amssymb,mathtools}
\usepackage[scr=rsfs,cal=euler]{mathalfa}
\usepackage{xfrac}
\usepackage[unicode,hidelinks,bookmarks=false]{hyperref}
\newcommand{\ie}{i.e.\ }
\newcommand{\cf}{cf.\ }
\newcommand{\resp}{respectively\ }
\newcommand{\Subsection}{\subsection}
\providecommand{\nobreakdash}{\nobreak\hbox{-}}
\setlength{\emergencystretch}{2em}
\multlinegap0pt
\usepackage{bm}
\usepackage{bbm}
\usepackage{dsfont}
\usepackage{xspace}
\usepackage{cancel}
\usepackage{mathtools}
\usepackage{amsthm}
\makeatletter
\@ifundefined{Theorem}{\newtheorem{Theorem}{Theorem}}{}
\@ifundefined{Proposition}{\newtheorem{Proposition}[Theorem]{Proposition}}{}
\makeatother
\makeatletter
\expandafter\ifx\csname Forget\endcsname\relax
\newenvironment{Forget}{}{}
\fi
\makeatother
\title[Stolz Positive Scalar Curvature Structure Groups, Proper Act]{Stolz Positive Scalar Curvature Structure Groups, Proper Actions and Equivariant 2-Types}
\author{Massimiliano Puglisi, Thomas Schick, Vito Felice Zenobi}
\date{Source: arXiv:2412.07955v3 (CC BY 4.0)}
\begin{document}
\maketitle

\noindent\emph{Source and licence.} Stolz Positive Scalar Curvature Structure Groups, Proper Actions and Equivariant 2-Types. Version arXiv:2412.07955v3 (CC BY 4.0), distributed under the Creative Commons Attribution 4.0 International licence, as recorded in the arXiv licence metadata for this exact version and shown on the abstract page https://arxiv.org/abs/2412.07955v3. Full rights notice: licenses/arxiv-2412.07955-CC-BY-4.0.txt.\par\smallskip

\noindent\textit{Editorial scope.} Counted unit: the Proposition whose source label is \texttt{universalspace} in the article (source0.tex lines 1071--1078, source label \texttt{universalspace}), delivered together with the article's own complete original proof of it (source0.tex lines 1081--1137, one \texttt{proof} environment of 57 lines). No mathematics is written, repaired or re-derived: statement and proof are the article's own text line for line. Required context retained uncounted: none. Every definition, display and label of those lines is retained unchanged. Omitted from this package and not redistributed: 1--1070, 1079--1080, 1138--1458. Nothing inside a retained excerpt is omitted or altered: every retained body is delivered exactly as the article's lines are written, so no omission receipt of any kind accompanies this package. Citations: the retained excerpts carry no citation command at all, so no bibliography entry of the article is transcribed and no reference is redistributed. Reference adaptation: none was needed and none was made; every reference inside the delivered unit resolves inside it, so no unresolved reference or citation is produced. Printed slips kept literal: no printed word, symbol, hypothesis or number of the counted statement or of its proof was altered. Layout and class: the article's own class is \texttt{sigma} with packages including ifpdf, xy; the wrapper is the lane's pinned \texttt{amsart} editorial wrapper at 10pt on A4, which restates the theorem-like environments the retained text needs and adds the article's own macro definitions that the retained text needs (none), with extra wrapper packages (bm, bbm, dsfont, xspace, cancel, mathtools). The delivered numbering is the unit's own. Assets: no retained span contains a figure environment, a graphics-inclusion command, a PDF-page inclusion, a table or a TikZ picture; no image file is copied into this package, the package declares no asset dependency and no image file is redistributed with it. Licence: version arXiv:2412.07955v3 of this article is distributed under the Creative Commons Attribution 4.0 International licence, as recorded in the arXiv licence metadata for this exact version and shown on the abstract page https://arxiv.org/abs/2412.07955v3. A full rights notice is delivered at licenses/arxiv-2412.07955-CC-BY-4.0.txt. Characters: every retained excerpt is pure ASCII, so the delivered engine, XeLaTeX with its default Latin Modern text font, reports no missing character.
	\begin{Proposition}\label{universalspace}
		For every $G$-CW-complex $Y$ and for every natural transformation $\Phi\colon
		\Pi_1(Y;\allowbreak G)\to \Pi$, there exists, unique up to $G$-equivariant
		homotopy, a $G$-equivariant cellular map
		$\varphi\colon Y\to B \Pi$
		such that
		$\varphi_{\#}=\Phi$.
	\end{Proposition}

	\begin{proof}
		We begin by defining the map $\varphi$ on the $G$-equivariant $0$-skeleton of $Y$ by putting
		$\varphi_{|Y_{(0)}}:=\smash{\Phi(\{e\})_{|Y_{(0)}}}$.
		Now, we proceed to define $\varphi$ on the $1$-skeleton of $Y$. For a $G$
		$1$-cell
		$c$ of the form $G/H\times [0,1]$ (with isotropy $H$) pick the ordinary
		$1$-cell $c_H:=H/H\times D^1$ contained in $Y^H$. Then $c_H$
		defines (if we choose an orientation of it) an element \smash{$\gamma \in
		\Pi_1\bigl(Y^H\bigr)_{|Y^H_{(0)}}$}. Define now \smash{$\phi\colon c_H\to \bigl(B\Pi^H\bigr)^{(1)}$} such that
		this map represents the element $\Phi(H)(\gamma)$ in
		$\Pi_1\bigl(B\Pi^H\bigr)|_{ (B\Pi^{ (0)} )^H}$. Since $H$ acts trivially on the target, this extends
		uniquely to a $G$-map on the $G$-cell $c$ with values in the
		$1$-skeleton of $B\Pi$. Of course, we could have picked a different base cell
		in the $G$-cell $c$. But because $\Phi$ is a natural transformation, the
		resulting map is independent of this ---up to the choice of the representative in the
		homotopy class of $\Phi(H)(\gamma)$ which had already to be made anyway.
		This defines the map $\phi$ on the $1$-skeleton of $Y$.
		
		To extend $\phi$ to $Y^{(2)}$, let $c$ be a $G$-$2$-cell of the form
		$G/H\times D^2$. Pick the corresponding ordinary $2$-cell $H/H\times
		D^2$ (with isotropy $H$)
		which is contained in $Y^{H}$. Its attaching map~${\psi\colon S^1\to Y^{H}}$
		is obviously contractible in $Y^{H}$ and hence trivial (i.e., a unit) in
		the fundamental groupoid~\smash{$\Pi_1\bigl(B\Pi^H\bigr)|_{ (B\Pi^{ (0)})^H}$} (we conjugate with a path in
		$Y^H$ to the $0$-skeleton to make~$\psi$
		represent an element of \smash{$\Pi_1\bigl(Y^H\bigr)_{B\Pi^H_{(0)}}$}, the triviality does not
		depend on the choice of such a path). Consequently, $\Phi(H)$ being a map of
		groupoids, also the image under $\Phi(H)$ of this element of the fundamental
		groupoid is trivial (a unit). By construction of $\phi$ on the $1$-skeleton,
		this image element is represented by $\phi\circ\psi$ (up to the chosen
		conjugation with a path to the $0$-skeleton), which is hence contractible in
		$B\Pi^H$. Extend $\phi$ over $H/H\times D^2$ using this contraction and then
		extend it $G$-equivariantly to $c=G/H\times D^2$.
		
		Inductively, we then extend $\phi$ over the $k$-skeleta of $Y$. The extension
		property now follows because each attaching map has contractible image by the
		vanishing of all higher homotopy groups of all components of all fixed point
		sets of $B\Pi$ for the various subgroups of $H$.
		
		By the very construction of $\phi$ on the $1$-skeleton, we have
		$\varphi_{\#}=\Phi$, because the morphism sets of
		\smash{$\Pi(H)=\Pi_1\bigl(B\Pi^H\bigr)|_{ (B\Pi^{ (0)})^H}$} are
		generated by the $1$-cells.
		
		The last step of the proof concerns uniqueness. Choose therefore a
		$G$-equivariant map
		$\varphi'\colon Y \to B \Pi$ such that
		$\varphi'_{\#}=\Phi$. We have to show that $\varphi'$ is $G$-homotopic to
		$\varphi$.
		
		It is immediate from the definition that if $\varphi'_{\#}=\varphi_{\#}$, then
		their restrictions to $Y_{(0)}$ are equal. The construction of the desired
		$G$-homotopy is now done inductively over the skeleta and follows the pattern
		in the non-equivariant case and for the construction part, making use of the
		condition that $\phi_\#'=\phi_\#$ for the extension over the $1$-skeleton and
		of the vanishing of higher homotopy groups for the further extensions.
	\end{proof}

\end{document}
